% =====================================================================================
%  Chapter 1 — Introduction
%  Thesis: Comparison of Control Strategies for Interaction-Force Aware Multirotor Teams
%
%  DRAFT, 23 August 2026 (writing-track item W.9). Written last of the unblocked chapters,
%  deliberately: an introduction promises what the body delivers, and Chapters 2-5 now exist
%  to be checked against.
%
%  RULES FOR THIS CHAPTER.
%  1. No result is claimed. Nothing has flown. The contribution is stated as what the thesis
%     sets out to establish, not as what it found.
%  2. THIS FILE IS THE SOURCE OF TRUTH for every claim it makes about prior work. Verify a
%     claim against the paper itself, never against a summary or ledger in docs/.
%  3. The contribution list must match docs/19_Contribution_Statement.md, including its
%     "what this thesis does not claim" table. If they drift apart, THIS FILE WINS.
%
%  Citation baseline: 2026-08-27 (see docs/thesis/CITATION_BASELINE.md).
%  After that date: only NEW cites or EDITED prior-work sentences need re-checking.
% =====================================================================================

\chapter{Introduction}
\label{ch:intro}

% -------------------------------------------------------------------------------------
\section{Motivation}
\label{sec:intro-motivation}

A Crazyflie~2.0 measures \SI{9}{\centi\metre} rotor-to-rotor, and when the interaction between
vehicles is left unmodelled, planners have adopted conservative separations --- Shi et
al.~\cite{shi2020neuralswarm} cite \SI{60}{\centi\metre} from earlier planning work as one such
example. The margin is not arbitrary caution. Below it, the downwash of one vehicle acting on another is strong enough to destabilise
the lower one, and the standard remedy is to keep the vehicles far enough apart that the effect can
be ignored.

That margin is expensive. It sets the density at which a team can operate, and therefore the size
of the space a given number of vehicles requires. Applications that motivate multirotor teams in
the first place --- inspection in confined spaces, cooperative transport, aerial docking --- are
precisely those in which the space is not available. In some of them, close approach is the task
rather than a hazard to be avoided.

Removing the margin requires a controller that accounts for the interaction rather than avoiding
it, and that is difficult for a specific reason. The interaction is not a small perturbation
around a known model: it is strongly nonlinear in the relative configuration of the team, growing
sharply as vehicles close and decaying quickly with lateral offset, and it is markedly asymmetric,
since the vehicle above is affected far less than the one below. It is also hard to derive from
first principles in the regime that matters. Simplified models exist for part of the problem ---
the mean flow from a quadrotor's propellers in hover is well approximated by a turbulent jet
beyond about two and a half motor-to-motor distances below it~\cite{bauersfeld2024airflow} --- but
formation flight operates closer than that, in a near field for which those authors report no
closed-form model.

The field's response has been to obtain the interaction from data rather than from theory, and it
has worked. Reported results include a reduction in worst-case height error by a factor of two or
better, up to four~\cite{shi2020neuralswarm}, sixteen vehicles flying at a minimum vertical
separation of
\SI{24}{\centi\metre}~\cite{shi2022neuralswarm2}, and average reductions of \SI{56}{\percent} in
vertical and \SI{36}{\percent} in three-dimensional tracking error for a pair of
vehicles~\cite{smith2023so2equivariant}. The margin has been substantially reduced.

% -------------------------------------------------------------------------------------
\section{Problem}
\label{sec:intro-problem}

A controller can obtain information about the interaction force in exactly two ways, and the
distinction organises the entire field.

It can \textbf{measure} the residual online, by comparing the acceleration the vehicle actually
experienced against the acceleration its nominal model predicts for the actuation it applied. This
requires no prior knowledge of the interaction and works equally against disturbances never seen
before. But the estimate exists only after the disturbance has begun to act, and in its established
form it depends on rotor-speed sensing.

Or it can \textbf{predict} the residual before it acts, from the relative states of its neighbours,
using a model learned offline. This can compensate ahead of the disturbance rather than after it,
but the model is only valid over the configurations its training data covered, and outside them its
behaviour is not merely degraded but unconstrained.

These are complementary rather than competing: one is reactive and general, the other predictive
and specific. A hybrid is therefore an obvious construction, and has been built --- for a single
vehicle, where its advantage over the reactive method alone was described as small without a
payload, and as leaving performance similar to the reactive method with
one~\cite{cobobriesewitz2026lindi}.

The practical question a system designer faces is which to use, and it has no answer in the
literature. Not because the individual results are weak, and not because the field avoids
systematic comparison --- recent work benchmarks online-learning estimators on a single
vehicle~\cite{gu2025comparative} and fractional-order and classical regulators for close-proximity
swarms~\cite{abro2025fractional}. It is that the results for \emph{these} families cannot be
composed: each is reported against its own baseline, on its own airframe, at its own separations,
on its own trajectories, under its own metrics.

% -------------------------------------------------------------------------------------
\section{Gap}
\label{sec:intro-gap}

\emph{Representative methods exist for reactive compensation, for predictive compensation, and for
their hybrid, and each is individually validated --- but, to the best of our knowledge, they have
not been evaluated against one another under identical conditions, and the field's reporting conventions make such a comparison
impossible to assemble after the fact.} A \SI{36}{\percent} improvement measured on two vehicles in
close proximity and a \SI{31}{\percent} improvement measured in a tight formation are not the same
quantity; neither can be subtracted from the other, and their ordering says nothing about which
method would prevail if both were flown under the same conditions.

This is not a criticism of any individual work. Each is internally valid, and none violated a
convention of the field. It is a property of the literature taken as a whole: the results were not
designed to compose, and establishing a ranking requires an experiment built for that purpose.

Two narrower gaps follow. The hybrid has been demonstrated only where the residual is
\emph{self-induced} --- a single vehicle, with effects such as blade flapping, drag and ground
forces, optionally carrying a slung payload --- and not
where it is produced by other vehicles and is therefore a function of a measurable relative state,
which is where the predictive component has the most to contribute. And the summation that
predictive architectures use to aggregate per-neighbour contributions is known to be an
approximation~\cite{shi2020neuralswarm,gielis2023aggregate}, yet the size of the penalty it carries
for a deployed model is not reported.

% -------------------------------------------------------------------------------------
\section{Research questions}
\label{sec:intro-rq}

\begin{description}
  \item[RQ1] Under identical hardware, trajectories, formations and separations, how do reactive,
    predictive and hybrid interaction-force compensation strategies compare in trajectory-tracking
    accuracy and residual-force rejection?

  \item[RQ2] Does the relative advantage of reactive versus predictive compensation depend on the
    interaction regime --- separation, relative velocity, and whether the disturbance is
    quasi-static or transient?

  \item[RQ3] How much prediction accuracy does a deep-sets residual model trained only on
    two-vehicle data lose when applied to three-vehicle formations, and is that loss small enough
    for the summed per-neighbour architecture to remain the right engineering choice? \emph{This
    is a measurement of the transfer penalty, not a test of whether superposition holds: it does
    not, and that is already established}~\cite{shi2020neuralswarm,gielis2023aggregate}.

  \item[RQ4] What are the sensing, computational and data requirements of each strategy, and is the
    added complexity of the hybrid justified in the multi-robot case?
\end{description}

RQ1 is the primary question and the one the experimental design serves. RQ2 carries the most
engineering value, because it converts a ranking into a rule about when to use which. RQ3 is
deliberately narrower than asking whether interactions superpose --- they do not, and that is
already reported --- and asks instead for the magnitude of the resulting error on a deployed model.
RQ4 exists because a method that wins by a small margin while requiring additional sensing, an
offline training campaign and a larger compute budget may still be the wrong choice.

The questions are stated with falsifiable hypotheses in Chapter~\ref{ch:related} and
Chapter~\ref{ch:setup}; each is answerable by the experiments described in
Chapter~\ref{ch:setup}, and a negative answer to any of them is reported as a result rather
than as a failure.

% -------------------------------------------------------------------------------------
\section{Contributions}
\label{sec:intro-contrib}

This thesis sets out to establish the following. Nothing here is a claim about results; the
experiments are described in Chapter~\ref{ch:setup} and reported in
the two- and multi-robot results chapters.

\begin{enumerate}
  \item \textbf{A controlled comparison.} To the best of our knowledge the first evaluation of
    reactive, predictive and hybrid compensation on the same airframe, the same trajectories, the same formations and the same
    frozen gains, with a single definition of the residual and a measured noise floor. The
    methodological content is in the word \emph{controlled}. For Strategies 0, 1, 2 and 4 it is
    unusually strong: those four are not four controllers but four settings of one, differing only
    in which residual term is active, so a difference between them cannot be attributed to
    anything else. \emph{This does not extend to Strategies 3 and 6}, which are structurally
    different controllers --- a feedback-linearising law and a receding-horizon optimiser --- and
    for which ``identical conditions'' means identical trajectories, formations and hardware but
    not identical structure. The design provides: a frozen scenario library so every strategy
    meets identical excitation; gains fixed before data collection so a difference cannot
    be a tuning artefact; one residual definition logged identically under every controller,
    including those that do not use it; and null conditions in which no interaction is expected, so
    that a positive result is separable from noise.

  \item \textbf{The first multi-robot evaluation of the reactive--predictive hybrid,} extending a
    result established for a single vehicle into the regime where its premise is strongest.

  \item \textbf{A quantification of what the summed per-neighbour approximation costs} when a
    model trained on two vehicles meets three. This measures a penalty whose \emph{existence} is
    already established in the literature~\cite{shi2020neuralswarm,gielis2023aggregate}; the
    contribution is its magnitude for a deployed model, not the finding that it exists.

  \item \textbf{An engineering account of cost}: sensing prerequisites, onboard computation and
    training-data requirements, reported alongside performance so that the comparison supports a
    decision rather than only a ranking.

\end{enumerate}

Supporting these, but explicitly \emph{not} a scientific contribution, is the apparatus that makes
them feasible: the frozen scenario library, the residual measurement path, an onboard network with
a weight-upload protocol, an offline training pipeline verified against the compiled controller,
and a simulation environment running the same controller source that flies. It is described in
Chapter~\ref{ch:setup} and offered for reuse. A thesis whose contribution is its tooling has
not answered a question, and it is listed here only so that the engineering effort is visible.

\paragraph{What this thesis does not claim.} No new control method is proposed: every strategy
follows published work, and the contribution is the comparison. The residual model is not offered
as an improvement on published architectures. Results concern a single airframe and a single class
of disturbance; generalisation beyond either is discussed rather than demonstrated. No claim rests
on simulation, which is used to disprove implementations and never to validate a method. And no
strategy is claimed to be best in general --- the useful outcome is conditional, which is RQ2.

\paragraph{A negative result would also be a contribution.} If the strategies do not separate, the
conclusion is that the choice of compensation family matters less than whether compensation is
present at all, and the engineering question becomes one of cost. The experimental design is built
so that this outcome is distinguishable from a failed experiment, which is what the null conditions
and the frozen gains are for.

\todo[color=green!25]{\textbf{Provenance, not pending work.} ``To the best of our knowledge'' is supported by two audits: an arXiv keyword
sweep and a cross-database sweep including a citation-graph scan. Both found adjacent comparisons, cited in §\ref{sec:rw-existing-comparisons},
but none spanning these families. Neither audit is exhaustive.}

% -------------------------------------------------------------------------------------
\section{Scope}
\label{sec:intro-scope}

The scope is deliberately narrow, because the contribution is a controlled comparison and every
degree of freedom left open is an alternative explanation for an observed difference.

\begin{table}[h]
  \centering
  \caption{Scope of this thesis.}
  \label{tab:intro-scope}
  \small
  \begin{tabular}{@{}p{0.46\linewidth}p{0.46\linewidth}@{}}
    \toprule
    In scope & Out of scope \\
    \midrule
    Inter-vehicle aerodynamic interaction, principally downwash & Wind, gusts, and other external disturbances \\
    Teams of two and three vehicles & Large swarms \\
    Homogeneous vehicles & Heterogeneous teams \\
    Tracking of pre-planned, collision-free formations & Interaction-aware planning; collision avoidance \\
    Onboard inference, offline training & Online or on-board learning \\
    The force residual & The moment residual, which is measured but not separately compensated \\
    \bottomrule
  \end{tabular}
\end{table}

% -------------------------------------------------------------------------------------
\section{Outline}
\label{sec:intro-outline}

\textbf{Chapter~\ref{ch:related}} surveys the field as an argument rather than a catalogue:
what the disturbance is, the two ways of obtaining information about it, why the injection point
into the control law matters, and why the individual results do not compose into a comparison.

\textbf{Chapter~\ref{ch:model}} defines the residual wrench that every strategy estimates, gives
both routes to it in one notation, and states the assumptions the formulation makes together with
where each is expected to fail.

\textbf{Chapter~\ref{ch:control}} describes the seven strategies and the point at which each
injects its knowledge of the interaction, separating in each case what is implemented here from the
published method it follows.

\textbf{Chapter~\ref{ch:setup}} describes the platform, the frozen scenario library, the
measurement chain and the protocol, and states the threats to validity the design is built against
--- including those it does not eliminate.

\textbf{the two- and multi-robot results chapters} report the two- and
three-vehicle results. \textbf{the discussion chapter} interprets them against the
research questions and the stated assumptions, and \textbf{the concluding chapter}
concludes.
